% 缩写与符号说明：论文中主要英文缩写及中英文全称、符号含义对照。
% 按缩略词字母顺序排列。正文中使用 \gls{xxx} 或 \acrlong{xxx} 等引用。

\newacronym[description=应用专用集成电路]{asic}{ASIC}{Application-Specific Integrated Circuit}
\newacronym[description=卷积神经网络]{cnn}{CNN}{Convolutional Neural Network}
\newacronym[description=直接内存访问]{dma}{DMA}{Direct Memory Access}
\newacronym[description=动态随机存取存储器]{dram}{DRAM}{Dynamic Random Access Memory}
\newacronym[description=设计空间探索]{dse}{DSE}{Design Space Exploration}
\newacronym[description=深度神经网络]{dnn}{DNN}{Deep Neural Network}
\newacronym[description=现场可编程门阵列]{fpga}{FPGA}{Field-Programmable Gate Array}
\newacronym[description=先就绪先服务（内存调度策略）]{frfcfs}{FR-FCFS}{First-Ready First-Come-First-Served}
\newacronym[description=通用矩阵乘]{gemm}{GEMM}{General Matrix Multiply}
\newacronym[description=中间表示]{ir}{IR}{Intermediate Representation}
\newacronym[description=一种轻量级数据交换格式]{json}{JSON}{JavaScript Object Notation}
\newacronym[description=键值缓存（自注意力中的缓存）]{kvcache}{KV Cache}{Key-Value Cache}
\newacronym[description=大语言模型]{llm}{LLM}{Large Language Model}
\newacronym[description=乘加运算]{mac}{MAC}{Multiply-Accumulate}
\newacronym[description=多层中间表示]{mlir}{MLIR}{Multi-Level Intermediate Representation}
\newacronym[description=片上网络]{noc}{NoC}{Network-on-Chip}
\newacronym[description=神经网络处理器]{npu}{NPU}{Neural Processing Unit}
\newacronym[description=处理单元]{pe}{PE}{Processing Element}
\newacronym[description=寄存器传输级]{rtl}{RTL}{Register Transfer Level}
\newacronym[description=静态随机存取存储器]{sram}{SRAM}{Static Random Access Memory}
\newacronym[description=张量处理单元]{tpu}{TPU}{Tensor Processing Unit}
\newacronym[description=虚拟通道]{vc}{VC}{Virtual Channel}
\newacronym[description=流控单元；NoC 中最小传输单位]{flit}{Flit}{Flow control digit}
\newacronym[description=输入特征图]{ifmap}{ifmap}{input feature map}
\newacronym[description=输出特征图]{ofmap}{ofmap}{output feature map}

% 以下为论文中出现的工具/平台名称，列出全称与中文说明
% \newacronym[description=周期级片上网络模拟器]{booksim2}{BookSim2}{BookSim2 NoC Simulator}
% \newacronym[description=多核映射与调度探索框架（HPCA 2024）]{gemini}{Gemini}{Gemini Mapping Framework}
% \newacronym[description=DRAM 时序与调度模拟器]{ramulator2}{Ramulator2}{Ramulator2 DRAM Simulator}
% \newacronym[description=Synopsys 硬件模拟平台]{zebu}{ZeBu}{ZeBu Hardware Emulation Platform}
% \newacronym[description=层间调度空间定义与探索（ISCA 2023）]{set}{SET}{Inter-layer Scheduling Space Definition and Exploration}
